\section*{Capítulo \ref{ch:b2:magmatic-series} --- Series magmáticas y contextos tectónicos}

\begin{solution}{exo:b2:magmatic-series:1}
(a) $\ce{SiO2} = 49.60$, álcalis $2.22 + 0.51 = 2.73$: entre 45 y 52 y por debajo de 5,
\emph{basalto} (B). (b) 58.5 y $3.9 + 1.6 = 5.5$: entre 57 y 63; la frontera
superior va allí de (57, 5.9) a (63, 7.0), y está en $5.9 + 1.1\times1.5/6 = 6.18$
para 58.5; 5.5 queda por debajo: \emph{andesita} (O2). (c) 72.0 y $4.0 + 4.2 = 8.2$:
más allá de la vertical en 69 por encima de 8 y por encima de la recta de (69, 8) a (77, 0), que
está en 5.0 para 72: \emph{riolita} (R). (d) 47.0 y $4.5 + 2.5 = 7.0$: en 47 el campo U1
está entre la recta de (45, 5) a (49.4, 7.3), en $5 + 2.3\times2/4.4 = 6.05$, y
la recta de (45, 9.4) a (49.4, 7.3), en $9.4 - 2.1\times2/4.4 = 8.45$. La roca está
en U1: una \emph{tefrita} si su norma contiene menos de un \qty{10}{\percent} de olivino, una
\emph{basanita} si contiene más.
\end{solution}

\begin{solution}{exo:b2:magmatic-series:2}
(a) Los óxidos suman 98.40 y el total con volátiles 100.00;
\ce{H2O} $1.40 < 2$ y \ce{CO2} $0.20 < 0.5$: el análisis es aceptable. (b) Factor
$100/98.40 = 1.01626$: \ce{SiO2} 52.44, \ce{TiO2} 1.22, \ce{Al2O3} 16.36,
\ce{FeO_T} 9.96, \ce{MgO} 6.71, \ce{CaO} 8.33, \ce{Na2O} 3.56, \ce{K2O} 1.42
(suma 100.00); álcalis $3.56 + 1.42 = 4.98$. (c) Antes: 51.60 y 4.90, justo por debajo de la
frontera basalto--andesita basáltica en 52: \emph{basalto}. Después: 52.44 y 4.98; en
52.44 la frontera superior de O1, de (52, 5) a (57, 5.9), está en $5 + 0.9\times0.44/5
= 5.08$: \emph{andesita basáltica} (O1). El recálculo ha hecho cruzar a la roca una
frontera: los nombres cerca de una línea dependen de él. (d) En 52.44 la recta de bajo K está en $0.3 +
0.045\times4.44 = 0.50$ y la de alto K en $1.2 + 0.095\times4.44 = 1.62$; 1.42
queda entre ambas: \emph{media en K}.
\end{solution}

\begin{solution}{exo:b2:magmatic-series:3}
(a) $Q + A + P = 90$: $Q = 25/90 = \qty{27.8}{\percent}$ y $100\,P/(A + P) =
100\times35/65 = 53.8$: campo 3b, un \emph{monzogranito}, un granito. La biotita y la
hornblenda no entran en la suma QAPF. (b) $Q + A + P = 65$: $Q = 5/65 =
\qty{7.7}{\percent}$ (entre 5 y 20) y $100\times55/60 = 91.7$ (por encima de 90): campo
10*, diorita cuarcífera, gabro cuarcífero o anortosita cuarcífera. El \omterm{def:b2:magmatic-series:qapf}{diagrama QAPF} por sí solo no
elige entre ellas; lo hacen la plagioclasa y los minerales máficos (aquí un
\qty{35}{\percent} de clinopiroxeno).
\end{solution}

\begin{solution}{exo:b2:magmatic-series:4}
La recta de bajo K es $0.3 + 0.045(\ce{SiO2} - 48)$, la de alto K $1.2 +
0.095(\ce{SiO2} - 48)$. (a) En 50: 0.39 y 1.39; $0.40$ está justo por encima de la recta inferior:
\emph{media en K}. (b) En 56: 0.66 y 1.96: \emph{media en K}. (c) En 62: 0.93 y 2.53;
$2.90 > 2.53$: \emph{alta en K}. (d) En 66: 1.11 y 2.91; $1.00 < 1.11$: \emph{baja en K}.
\end{solution}

\begin{solution}{exo:b2:magmatic-series:5}
(a) Basalto: $A + F + M = 3.2 + 14.5 + 4.5 = 22.2$; $A = 14.4$, $F = 65.3$, $M = 20.3$.
Andesita: $5.05 + 6.1 + 3.8 = 14.95$; $A = 33.8$, $F = 40.8$, $M = 25.4$. (b) La línea
nunca pasa de $F = 54.5$; el basalto, con $F = 65.3$, cae del lado de F
sea cual sea su $A$: \emph{toleítico}, un ferrobasalto de una tendencia de Fenner. Para la
andesita, en $A = 33.8$ la línea está entre (49.6, 35.0) y (52.5, 30.0), en $F =
49.6 + 2.9\times1.2/5 = 50.3$; $40.8 < 50.3$: \emph{calcoalcalina}. (c) La magnetita,
que retira hierro férrico de un \omterm{def:b2:magma-emplacement:magma}{magma} oxidado desde el principio e impide que $F/M$
aumente tan deprisa como en una toleita.
\end{solution}

\begin{solution}{exo:b2:magmatic-series:6}
MORB: $\mathrm{Nb_N} = 4/0.658 = 6.08$, $\mathrm{La_N} = 4.12/0.648 = 6.36$, razón
0.956. \omterm{def:b2:elements-in-the-earth:morb-oib}{Basalto de isla oceánica}: $38/0.658 = 57.8$, $30.8/0.648 = 47.5$, razón 1.215. Basalto
de arco: $3/0.658 = 4.56$, $6.1/0.648 = 9.41$, razón 0.484. Solo el basalto de arco tiene una
anomalía de niobio, un valle en Nb en un diagrama de araña: el niobio está empobrecido respecto de
su vecino La: el \omterm{def:b2:magma-emplacement:magma}{magma} de arco recibió elementos de la placa, pero poco
niobio. El MORB es casi plano; el \omterm{def:b2:elements-in-the-earth:morb-oib}{basalto de isla oceánica} es más rico
en Nb que en La.
\end{solution}

\begin{solution}{exo:b2:magmatic-series:7}
(a) $\mathrm{Ti} = 1.6\times0.5993\times10^4 = \qty{9589}{ppm}$; Ti/100 = 95.9, Zr
110, $3\mathrm Y = 96$; suma 301.9; coordenadas 31.8, 36.4, 31.8: campo B. (b) Ti
16\,780, luego 167.8, 220, 87; suma 474.8; 35.3, 46.3, 18.3: campo D. (c) Ti 4794, luego
47.9, 50, 60; suma 157.9; 30.3, 31.7, 38.0: campo B. El basalto (b) es un basalto intraplaca
(de una isla oceánica o de un continente); (a) y (c), en el campo B, pueden ser MORB, toleitas
de arco o basaltos calcoalcalinos: el diagrama no puede decidir, y hace falta otro (un
diagrama de araña, la anomalía de Nb).
\end{solution}

\begin{solution}{exo:b2:magmatic-series:8}
(a) MORB $142/164 = 0.866$, \omterm{def:b2:elements-in-the-earth:morb-oib}{basaltos de isla oceánica} $4/173 = 0.023$, basaltos de arco $101/201 =
0.502$. (b) Con probabilidades a priori iguales, la regla de Bayes da $0.866/(0.866 + 0.023 + 0.502) =
0.866/1.391 = 0.623$. (c) Que es muy improbable que el basalto sea un \omterm{def:b2:elements-in-the-earth:morb-oib}{basalto de isla oceánica}
(probabilidad $0.023/1.391 = 0.017$), pero no si es un MORB o un basalto de arco:
aproximadamente tres posibilidades de cinco frente a dos.
\end{solution}

\begin{solution}{exo:b2:magmatic-series:9}
Moles por cada \qty{100}{g}: \ce{Na2O} $5.5/61.98 = 0.08874$, \ce{K2O} $4.5/94.20 =
0.04777$, \ce{Al2O3} $10.5/101.96 = 0.10298$. (a) $(0.08874 + 0.04777)/0.10298 =
0.13651/0.10298 = 1.326 > 1$: la riolita es \emph{peralcalina}. (b) El feldespato puede
tomar $0.10298$ de los $0.13651$ moles de óxido alcalino; el exceso, $0.03353$, es
$0.03353/0.13651 = 0.246$, una cuarta parte de los álcalis, que deben entrar en otros minerales.
\end{solution}

\begin{solution}{exo:b2:magmatic-series:10}
(a) $s = (\lambda_F - \lambda_M)/(\lambda_F - \lambda_A)$: para $\lambda_F = 0.5$, $(0.5
- 0.8)/(0.5 - 0.4) = -3.0$; para $0.7$, $(-0.1)/(0.3) = -0.333$. (b) Despejando de
$s(\lambda_F - 0.4) = \lambda_F - 0.8$ se obtiene $\lambda_F = (0.4s - 0.8)/(s - 1)$:
Islandia, $s = -1.10$: $(-0.44 - 0.8)/(-2.10) = 0.590$; Cascadas, $s = -0.72$:
$(-0.288 - 0.8)/(-1.72) = 0.633$. (c) El \omterm{def:b2:magma-emplacement:magma}{magma} de las Cascadas, con el mayor
$\lambda_F$, pierde su hierro más deprisa: la magnetita lo retira. (d) $\ln(F/M)$ crece al
ritmo $\lambda_M - \lambda_F$ y $\ln(A/F)$ al ritmo $\lambda_F - \lambda_A$; así, por
cada unidad que gana $\ln(A/F)$, $\ln(F/M)$ gana $-s$: 1.10 en Islandia, 0.72 en las
Cascadas. $F/M$ aumenta a lo largo de ambas tendencias, pero más deprisa en la toleítica, la tendencia
de Fenner.
\end{solution}

\begin{solution}{exo:b2:magmatic-series:11}
(a) BHVO-2, moles por cada \qty{100}{g}: \ce{SiO2} 0.82557, \ce{Al2O3} 0.13182, \ce{FeO}
0.15521, \ce{MgO} 0.18007, \ce{CaO} 0.20328, \ce{Na2O} 0.03580, \ce{K2O} 0.00545.
Ortoclasa 0.00545 y albita 0.03580; \ce{Al2O3} restante $0.13182 - 0.04125 = 0.09057$,
todo en anortita con otro tanto de \ce{CaO}; \ce{CaO} restante $0.11271$ en diópsido;
$(\ce{Mg,Fe})\ce{O}$ restante $0.33528 - 0.11271 = 0.22257$ en hiperstena. Sílice
necesaria: $6\times0.00545 + 6\times0.03580 + 2\times0.09057 + 2\times0.11271 + 0.22257 =
0.87663$; disponible 0.82557; faltan 0.05106. Cada $(\ce{Mg,Fe})\ce{O}$ que pasa de
hiperstena a olivino ahorra media sílice, como mucho $0.22257/2 = 0.11129$: basta con convertir
$0.10212$, y $0.12045$ se queda en hiperstena. La norma contiene olivino
e hiperstena, y ni cuarzo ni nefelina: la roca está saturada en sílice en el
sentido de la definición, una \emph{toleita olivínica} con hiperstena normativa, de la
\omterm{def:b2:magmatic-series:series}{serie toleítica}, como las que emite un volcán hawaiano en su etapa de escudo. (b) Moles: 0.75732, 0.14712, 0.16008, 0.21092, 0.17832, 0.06131,
0.01592. Or 0.01592, ab 0.06131, an 0.06988, di 0.10843, hy $0.37100 - 0.10843 =
0.26257$; sílice necesaria $0.09552 + 0.36786 + 0.13976 + 0.21686 + 0.26257 = 1.08257$;
disponible 0.75732; faltan 0.32525. Convertir toda la hiperstena en olivino ahorra
0.13129, y quedan 0.19396; cada mol de \ce{Na2O} que pasa de albita a nefelina
libera 4 \ce{SiO2}, así que $0.19396/4 = 0.04849$
de los 0.06131 moles de \ce{Na2O} de la albita pasan a nefelina. Nefelina en la norma: la
roca está subsaturada en sílice, un \emph{basalto alcalino}, de la \omterm{def:b2:magmatic-series:series}{serie alcalina}.
\end{solution}

\begin{solution}{exo:b2:magmatic-series:12}
(a) Basaltos de arco $116/141 = 0.823$, MORB $21/101 = 0.208$, \omterm{def:b2:elements-in-the-earth:morb-oib}{basaltos de isla oceánica}
$3/137 = 0.022$. (b) $0.823/(0.823 + 0.208 + 0.022) = 0.823/1.053 = 0.782$.
(c) Numerador $0.1\times0.823 = 0.0823$; denominador $0.0823 + 0.6\times0.208 +
0.3\times0.022 = 0.0823 + 0.1248 + 0.0066 = 0.2137$; probabilidad $0.385$. (d) La
misma prueba sobre la misma roca da 0.78 o 0.39 según lo común que sea cada contexto
donde se encontró la roca: un \omterm{def:b2:magmatic-series:discrimination}{diagrama de discriminación} da una verosimilitud, y el
contexto geológico debe aportar el resto.
\end{solution}

\begin{solution}{pb:b2:magmatic-series:1}
\noindent\textbf{Parte I.}

\textbf{1.} Óxidos: lava 1 99.17, lava 2 98.65, lava 3 98.20; con los volátiles
99.97, 99.80 y 99.45, todos cerca de 100. El agua (0.70, 1.10, 1.20) está por debajo del
\qty{2}{\percent} y el dióxido de carbono (0.10, 0.05, 0.05) por debajo del \qty{0.5}{\percent}: las
tres son aceptables.

\textbf{2.} Factor $100/99.17 = 1.00837$: \ce{SiO2} 50.82, \ce{Na2O} 2.92, \ce{K2O}
0.63, álcalis 3.55.

\textbf{3.} Lava 2: $100/98.65 = 1.01368$, \ce{SiO2} 59.40, \ce{Na2O} 3.95, \ce{K2O}
1.17, álcalis 5.12. Lava 3: $100/98.20 = 1.01833$, \ce{SiO2} 66.50, \ce{Na2O} 4.48,
\ce{K2O} 1.88, álcalis 6.36.

\textbf{4.} Lava 1 (50.82, 3.55): basalto. Lava 2 (59.40, 5.12; la frontera superior de O2
está allí en $5.9 + 1.1\times2.40/6 = 6.34$): andesita. Lava 3 (66.50, 6.36; entre 63
y la recta de (69, 8) a (77, 0), por debajo de la recta de (63, 7) a (69, 8)): dacita.

\textbf{5.} Sin recalcular: (50.4, 3.52), (58.6, 5.05), (65.3, 6.25): los mismos campos. El
recálculo desplaza cada punto como mucho un \qty{1.2}{\percent} de sílice y no cambia
nada aquí, a diferencia del \cref{exo:b2:magmatic-series:2}.

\textbf{6.} El \omterm{def:b2:magmatic-series:qapf}{diagrama QAPF}, a partir de la moda medida. Una roca plutónica de composición
andesítica es rica en plagioclasa, con poco feldespato alcalino o cuarzo: debería
caer en el campo 10 o 10*, una diorita o una diorita cuarcífera (decide su moda).

\noindent\textbf{Parte II.}

\textbf{7.} Lava 1: rectas en 0.43 y 1.47, $\ce{K2O} = 0.63$: media en K. Lava 2: 0.81
y 2.28, 1.17: media en K. Lava 3: 1.13 y 2.96, 1.88: media en K.

\textbf{8.} \ce{FeO_T} y \ce{MgO} sin volátiles: 8.67 y 7.97; 6.18 y 3.85; 4.07
y 1.63. Lava 1: $3.55 + 8.67 + 7.97 = 20.19$; $A = 17.6$, $F = 42.9$, $M = 39.5$.
Lava 2: 15.15; 33.8, 40.8, 25.4. Lava 3: 12.06; 52.7, 33.7, 13.5.

\textbf{9.} Lava 1, $A = 17.6$: la línea entre (42.8, 17.1) y (46.9, 18.1) está en
$F = 42.8 + 4.1\times0.5/1.0 = 44.9$; $42.9 < 44.9$, lado calcoalcalino, cerca de la
línea. Lava 2, $A = 33.8$: $F = 49.6 + 2.9\times1.2/5 = 50.3$; $40.8 < 50.3$,
calcoalcalina. Lava 3, $A = 52.7$: $F = 40.8 - 2.9\times2.7/5 = 39.2$; $33.7 < 39.2$,
calcoalcalina.

\textbf{10.} $\ce{FeO_T}/\ce{MgO}$: $8.67/7.97 = 1.09$, $6.18/3.85 = 1.61$, $4.07/1.63 =
2.50$. La razón aumenta, como en todo basalto que se diferencia, pero el propio hierro
baja de 8.67 a 4.07, y $F$ baja en el triángulo de 42.9 a 33.7: no hay enriquecimiento
en hierro. Una \omterm{def:b2:magmatic-series:series}{serie toleítica} subiría primero hacia F, como el ferrobasalto del
\cref{exo:b2:magmatic-series:5}.

\textbf{11.} La \omterm{def:b2:magmatic-series:series}{serie calcoalcalina}. La magnetita, al cristalizar de un \omterm{def:b2:magma-emplacement:magma}{magma} oxidado,
retiró hierro junto con el magnesio; las especies oxidadas proceden de la placa que se deshidrata en una
zona de subducción.

\noindent\textbf{Parte III.}

\textbf{12.} $\mathrm{Nb_N} = 2.0/0.658 = 3.04$, $\mathrm{La_N} = 7/0.648 = 10.80$,
$\mathrm{Th_N} = 1.1/0.0795 = 13.84$.

\textbf{13.} $3.04/10.80 = 0.281$, muy por debajo de las medianas de los MORB (0.956) y de los
\omterm{def:b2:elements-in-the-earth:morb-oib}{basaltos de isla oceánica} (1.215), e incluso por debajo del basalto de arco mediano (0.484): un profundo
valle de niobio.

\textbf{14.} $13.84/3.04 = 4.55$; el MORB mediano del diagrama de araña tiene $2.64/6.08 =
0.43$. El torio, como los demás elementos enriquecidos en los basaltos de arco, lo añadió la
placa que subduce.

\textbf{15.} \ce{TiO2} sin volátiles $0.95\times1.00837 = 0.958$; $\mathrm{Ti} =
0.958\times5993 = \qty{5741}{ppm}$; Ti/100 = 57.4, Zr 70, $3\mathrm Y = 54$, suma 181.4:
31.6, 38.6, 29.8, en el campo B, compartido por los MORB, las toleitas de arco y los basaltos calcoalcalinos.
Excluye un basalto intraplaca pero no decide entre una dorsal y un arco. El
diagrama se construyó para basaltos: en una andesita y una dacita, los óxidos han retirado el titanio
y la cristalización ha concentrado el circonio, y sus campos (la lava 2
caería en C, la lava 3 fuera de todo campo) no significarían nada.

\textbf{16.} Una \omterm{def:b2:magmatic-series:series}{serie calcoalcalina}, media en K, con un profundo valle de niobio y una
$\mathrm{Th_N}/\mathrm{Nb_N}$ alta, en un volcán que emitió de basalto a dacita: el volcán
creció en un arco volcánico sobre una zona de subducción.

\noindent\textbf{Parte IV.}

\textbf{17.} $C_L = C_0/F$ para un elemento perfectamente incompatible: $F = 0.63/1.17 =
0.538$.

\textbf{18.} Rubidio $10/20 = 0.500$; torio $1.1/2.1 = 0.524$. (Los \omterm{def:b2:elements-in-the-earth:trace}{elementos en trazas},
dados tal como se analizaron, cambiarían menos de un \qty{1}{\percent} con el
recálculo.)

\textbf{19.} $C_L/C_0 = F^{D - 1}$: $125/70 = 1.786 = F^{-0.9}$, $F = 1.786^{-1/0.9} =
0.525$.

\textbf{20.} A partir de la potasa, $F = 1.17/1.88 = 0.622$ de la lava 2 a la lava 3. Estroncio:
$380/480 = 0.792 = 0.622^{D - 1}$, luego $D - 1 = \ln 0.792/\ln 0.622 = -0.234/-0.474 =
0.49$ y $D \approx 1.49$: el estroncio se ha vuelto compatible. La plagioclasa, que capta
estroncio, cristalizaba en abundancia.

\textbf{21.} A partir de la potasa, $F = 0.63/1.88 = 0.335$ (rubidio $10/32 = 0.31$, torio
$1.1/3.4 = 0.32$): cerca de dos tercios del basalto, $1 - 0.335 = 0.665$, cristalizaron
antes de la dacita.

\textbf{22.} La potasa asimilada elevaría el \ce{K2O} de la lava 2 por encima de lo que da
la cristalización sola, de modo que $F = 0.63/1.17$ saldría demasiado bajo, y la masa de
cristales demasiado alta.

\textbf{23.} $(0.538 + 0.500 + 0.524 + 0.525)/4 = 2.087/4 = 0.522$; cristales $1 - 0.522
= 0.478$ del basalto, es decir, \textbf{\qty[text-series-to-math=true]{48}{kg} de
cristales por cada \qty{100}{kg} de basalto} para formar la andesita.
\end{solution}
